\chapter{Descrizione dello stage}
\label{chap:descrizione-stage}

\section{Descrizione del contesto aziendale}

Nonostante Spazio Dev sia una realtà di recente fondazione, l'organizzazione interna è caratterizzata da processi ben definiti e da una forte attenzione alla qualità del lavoro svolto e alla soddisfazione del cliente.

Durante il periodo di stage ho collaborato con un team composto da tre persone. In un contesto di dimensioni ridotte, la comunicazione costante e il coordinamento tra i membri del gruppo risultano fondamentali per rispondere efficacemente alle esigenze del cliente e garantire il successo del progetto.

\subsection{Modello Agile}

L'azienda adotta il framework Scrum come metodologia di sviluppo Agile, organizzando il lavoro in sprint della durata di una settimana.

All'inizio di ogni sprint viene svolta una riunione di pianificazione alla quale partecipa anche il cliente. Durante questo incontro vengono raccolti i nuovi requisiti, valutati gli avanzamenti del progetto e definite le attività da svolgere nella settimana successiva. I compiti vengono quindi distribuiti tra i membri del team in base alle competenze e alle priorità individuate.

Il ruolo di \textit{Product Owner} è svolto dai titolari dell'azienda, che fungono da riferimento per il cliente, raccolgono i requisiti e stabiliscono le priorità delle attività da realizzare.

Il ruolo di \textit{Scrum Master} è ricoperto da uno dei titolari, il quale ha il compito di facilitare la comunicazione all'interno del team, rimuovere eventuali ostacoli e assicurare il corretto svolgimento del processo Scrum.

Il \textit{Development Team} è composto da tutti i membri del gruppo di lavoro, che collaborano attivamente allo sviluppo del prodotto, condividendo competenze e responsabilità per raggiungere gli obiettivi prefissati entro ciascuno sprint.

\subsection{Issue Tracking, Version Control e Code Review}

Per la gestione delle attività di sviluppo, l'azienda utilizza Plane come sistema di \textit{issue tracking}. Le attività vengono organizzate attraverso board in stile Kanban e assegnate ai membri del team in base alle competenze e alle priorità definite durante le riunioni di pianificazione degli sprint.

I task sono inoltre raggruppati in moduli corrispondenti alle diverse aree del progetto, facilitando la suddivisione del lavoro e il monitoraggio dell'avanzamento delle varie componenti applicative. In accordo con il cliente, vengono inoltre stabilite specifiche milestone che scandiscono il rilascio progressivo delle funzionalità previste.

Per il controllo di versione del codice sorgente viene utilizzato Git, con repository ospitati sulla piattaforma Gitea. Le modifiche vengono integrate nel progetto attraverso un processo di \textit{code review} basato su \textit{pull request}, che prevede la verifica e l'approvazione da parte di almeno uno sviluppatore senior prima della fusione nel ramo di destinazione.

A supporto di questo processo viene inoltre impiegato un sistema automatizzato di revisione del codice basato su intelligenza artificiale, in grado di individuare potenziali problematiche, evidenziare possibili miglioramenti e contribuire al mantenimento della qualità del codice.

Il flusso di sviluppo adottato segue il modello \textit{GitFlow}, articolato nei seguenti rami principali:

\begin{itemize}
    \item \textbf{main}: contiene il codice stabile e pronto per il rilascio in produzione;
    \item \textbf{develop}: rappresenta il ramo di integrazione delle funzionalità in fase di sviluppo e costituisce la base per le future release;
    \item \textbf{feature}: rami temporanei dedicati allo sviluppo di nuove funzionalità o miglioramenti, creati a partire da \textit{develop} e successivamente reintegrati nello stesso;
    \item \textbf{release}: rami utilizzati per la preparazione di una nuova versione del software, nei quali vengono effettuate le attività di testing finale e correzione dei bug;
    \item \textbf{hotfix}: rami creati a partire da \textit{main} per la risoluzione rapida di problemi critici riscontrati in ambiente di produzione; una volta completati, vengono reintegrati sia in \textit{main} sia in \textit{develop}.
\end{itemize}

\subsection{Il solver}
Come anticipato il progetto proposto dal cliente è ambizioso e complesso, di fatti il team ci lavora da mesi. La parte che verrà trattata in questo stage è 
quella relativa allo sviluppo del solver di scheduling, ovvero il componente software che si occupa di risolvere i problemi di ottimizzazione, in 
particolare quelli legati alla pianificazione e all'allocazione delle risorse.
Questo modulo viene integrato tramite api con la libreria Drogon, che consente di esporre le funzionalità del solver attraverso un'interfaccia, 
l'accesso e l'utilizzo da parte del prodotto principale scritto in PHP. 
Nonostante ciò durante lo stage è stata sviluppata un interfaccia in qt di debug utile per testare e verificare il corretto funzionamento del solver.
DA FARE DA FARE

\section{Analisi preventiva dei rischi}

Durante la fase di analisi iniziale sono stati individuati alcuni possibili rischi a cui si potrà andare incontro.
Si è quindi proceduto a elaborare delle possibili soluzioni per far fronte a tali rischi.

\begin{risk}{Performance del simulatore hardware}
    \riskdescription{le performance del simulatore hardware e la comunicazione con questo potrebbero risultare lenti o non abbastanza buoni da causare il fallimento dei test}
    \risksolution{coinvolgimento del responsabile a capo del progetto relativo il simulatore hardware}
    \label{risk:hardware-simulator} 
\end{risk}

\section{Requisiti e obiettivi}

\begin{center}
    \rowcolors{1}{}{tableGray}
    \begin{longtable}{|p{2.25cm}|p{7.75cm}|p{2.25cm}|}
    \hline
    \multicolumn{1}{|c|}{\textbf{A}} & \multicolumn{1}{c|}{\textbf{B}}\\ 
    \hline 
    \endfirsthead
    \rowcolor{white}
    \multicolumn{3}{c}{{\bfseries \tablename\ \thetable{} -- Continuo della tabella}}\\
    \hline
    \multicolumn{1}{|c|}{\textbf{A}} & \multicolumn{1}{c|}{B}\\ \hline 
    \endhead
    \hline
    \rowcolor{white}
    \multicolumn{3}{|r|}{{Continua nella prossima pagina...}}\\
    \hline
    \endfoot
    \endlastfoot 
    
    AA & BB \\
    \hline
    AA & BB \\
    \hline
    AA & BB \\
    \hline
    AA & BB \\
    \hline
    \hiderowcolors
    \caption{Lorem.}
    \label{tab:requisiti_obbiettivi}
    \end{longtable}
\end{center}

\section{Pianificazione}
\begin{figure}[H]
    \centering
    \includegraphics[alt={Testo alternativo dell'immagine}, width=0.5\columnwidth]{img/pk_estate.jpeg}
    \caption{Caption}
    \label{fig:pk_estate_2}
\end{figure}
\lipsum[1]

\subsection{subsection}
\lipsum[1]

\subsubsection{subsubsection}
\lipsum[1]

\paragraph{paragraph}
\lipsum[1]

\newpage